\documentclass[a4paper,11pt]{article}

\usepackage{amsmath,amssymb}
\usepackage{mathpartir}
\usepackage{geometry}

\geometry{margin=25mm}

\title{Simply Typed Lambda Calculus}
\author{}
\date{}

\begin{document}

    \maketitle

    \section{Syntax}

    \subsection{Types}

    \[
        A,B ::= \iota \mid A \rightarrow B
    \]

    where

    \begin{itemize}
        \item $\iota$ is the base type.
        \item $A \rightarrow B$ is the function type.
    \end{itemize}

    \subsection{Terms}

    \[
        t ::= x
        \mid \lambda.t
        \mid t\;t
    \]

    Variables are represented using de Bruijn indices.

    \section{Typing Rules}

    Typing judgement

    \[
        \Gamma \vdash t : A
    \]

    \begin{mathpar}

        \inferrule*[right=T-Var]
        {\Gamma(x)=A}
        {\Gamma \vdash x : A}

        \inferrule*[right=T-Lam]
        {\Gamma,A \vdash t:B}
        {\Gamma \vdash \lambda.t : A \rightarrow B}

        \inferrule*[right=T-App]
        {\Gamma \vdash t : A \rightarrow B \\
        \Gamma \vdash u : A}
        {\Gamma \vdash t\,u : B}

    \end{mathpar}

    \section{Operational Semantics}

    Evaluation judgement

    \[
        t \longrightarrow t'
    \]

    Values

    \[
        v ::= \lambda.t
    \]

    \begin{mathpar}

        \inferrule*[right=E-AppAbs]
        {}
        {(\lambda.t)\;v \longrightarrow t[v]}

        \inferrule*[right=E-App1]
        {t_1 \longrightarrow t_1'}
        {t_1\;t_2 \longrightarrow t_1'\;t_2}

        \inferrule*[right=E-App2]
        {t_2 \longrightarrow t_2'}
        {v\;t_2 \longrightarrow v\;t_2'}

    \end{mathpar}

    \section{Metatheory}

    The Lean formalization proves the following results.

    \begin{itemize}
        \item Rename
        \item Substitution
        \item Progress
        \item Preservation
        \item Soundness
    \end{itemize}

    The final theorem is:

    \[
        \text{Well-typed programs do not get stuck.}
    \]

\end{document}